%!TEX root = ../hutbthesis_main.tex

% 设置中文摘要
\keywordscn{固定翼飞行器，着陆控制，人机交互强化}
\begin{abstractzh}

固定翼飞行器着陆阶段是飞行全过程中风险较高、控制要求较强的关键环节，易受到速度、下沉率、姿态偏差、侧偏误差及地形环境等因素影响。本文以基于人机交互强化的固定翼飞行器着陆控制系统为研究对象，围绕系统组成与需求分析、动力学模型建立、仿真平台实现、多工况结果分析以及交互强化策略设计等内容展开研究。首先，对着陆任务场景、关键参数及影响因素进行了分析；其次，建立了着陆阶段动力学模型和仿真模型，并完成基准工况校核；再次，通过平坦、起伏、山脊、谷地和阶跃等不同地形工况的仿真对比，分析了飞行器在着陆末段的离地高度、速度和姿态变化规律，识别出阶跃地形为典型不利工况；最后，提出了风险识别、状态提示、操纵辅助和约束保护相结合的人机交互强化策略。研究表明，该方法能够在一定程度上提高固定翼飞行器在复杂工况下的着陆安全性与稳定性，为后续控制系统优化提供参考。

\end{abstractzh}
